\section{The GFIC Diagnostic Protocol}
\label{sec:method}

\paragraph{Design principle.}
Each axis targets one link of the causal chain
\emph{perception $\to$ concept $\to$ action $\to$ robustness under domain shift $\to$
locatable failure}, is anchored to an independent literature (\cref{sec:related}), and is
required to pass a causal-validity check before being trusted, not merely to produce a
number that correlates with an outcome. All four axes are operationalized on a shared
stimulus set (\cref{sec:experiments}) so that cross-model comparisons are always
within-model-normalized, never raw-score comparisons.

\begin{figure*}[t]
\centering
\newcommand{\gficbox}[1]{\tikz[baseline=-0.5ex]\node[draw, rounded corners, align=center, inner sep=3pt, minimum height=6mm, font=\scriptsize]{#1};}
\newcommand{\gfichl}[1]{\tikz[baseline=-0.5ex]\node[draw=cvprblue, thick, fill=cvprblue!12, rounded corners, align=center, inner sep=3pt, minimum height=6mm, font=\scriptsize]{#1};}
\begin{minipage}[t]{0.48\textwidth}
\textbf{\small (a) G: representation vs.\ objective segmentation GT}\\[3pt]
\begin{tikzpicture}[every node/.style={font=\scriptsize}, box/.style={draw, rounded corners, align=center, minimum height=6mm, inner sep=3pt}, hl/.style={box, fill=cvprblue!12, draw=cvprblue, thick}, arr/.style={-{Latex[length=1.5mm]}, thick}]
\node[box] (frame) {Frame};
\node[hl, right=5mm of frame] (rep) {Repr.};
\node[box, above right=2mm and 5mm of rep] (probe) {Linear probe};
\node[box, below right=2mm and 5mm of rep] (sam) {SAM pseudo-GT};
\node[box, right=17mm of rep] (cmp) {mIoU compare};
\draw[arr] (frame) -- (rep);
\draw[arr] (rep) -- (probe);
\draw[arr] (frame.north) .. controls +(0,10mm) and +(-8mm,3mm) .. (sam.west);
\draw[arr] (probe) -- (cmp);
\draw[arr] (sam) -- (cmp);
\end{tikzpicture}\\[2pt]
{\scriptsize selectivity $=$ mIoU(trained) $-$ mIoU(rand-init)}
\end{minipage}%
\hfill
\begin{minipage}[t]{0.48\textwidth}
\textbf{\small (b) F: occlude the hazard entity, check necessity}\\[3pt]
\begin{tikzpicture}[every node/.style={font=\scriptsize}, box/.style={draw, rounded corners, align=center, minimum height=6mm, inner sep=3pt}, hl/.style={box, fill=cvprblue!12, draw=cvprblue, thick}, arr/.style={-{Latex[length=1.5mm]}, thick}]
\node[box] (ghost) {Ghost frame};
\node[hl, right=5mm of ghost] (occ) {Occlude entity};
\node[box, right=5mm of occ] (act) {Action};
\node[box, below=5mm of ghost] (clean) {Clean frame};
\draw[arr] (ghost) -- (occ);
\draw[arr] (occ) -- (act);
\draw[arr] (clean.east) -| (act.south);
\end{tikzpicture}\\[2pt]
{\scriptsize $R = (a_{occ}-a_{clean})/(a_{ghost}-a_{clean})$}
\end{minipage}

\vspace{6pt}
\begin{minipage}[t]{0.48\textwidth}
\textbf{\small (c) I: same scene, two renderings, two orderings}\\[3pt]
\begin{tikzpicture}[every node/.style={font=\scriptsize}, box/.style={draw, rounded corners, align=center, minimum height=6mm, inner sep=3pt}, hl/.style={box, fill=cvprblue!12, draw=cvprblue, thick}, arr/.style={-{Latex[length=1.5mm]}, thick}]
\node[box] (sim) {CARLA render};
\node[box, below=4mm of sim] (real) {Real-style render};
\node[hl, right=7mm of sim, yshift=-5mm] (model) {Model};
\node[box, right=9mm of model, yshift=5mm] (repD) {Repr.\ $D_L$};
\node[box, right=9mm of model, yshift=-5mm] (behD) {Behav.\ $\Delta v_{cmd}$};
\draw[arr] (sim) -- (model);
\draw[arr] (real) -- (model);
\draw[arr] (model) -- (repD);
\draw[arr] (model) -- (behD);
\end{tikzpicture}\\[2pt]
{\scriptsize the two orderings can disagree}
\end{minipage}%
\hfill
\begin{minipage}[t]{0.48\textwidth}
\textbf{\small (d) C: patch one layer, read the recovery profile}\\[3pt]
\begin{tikzpicture}[every node/.style={font=\scriptsize}, box/.style={draw, rounded corners, align=center, minimum height=6mm, minimum width=5mm, inner sep=2pt}, hl/.style={box, fill=cvprblue!12, draw=cvprblue, thick}, arr/.style={-{Latex[length=1.5mm]}, thick}]
\node[box] (l0) {L0};
\node[box, right=2mm of l0] (l1) {L1};
\node[hl, right=2mm of l1] (l2) {L2};
\node[box, right=2mm of l2] (l3) {L3};
\node[box, right=9mm of l3] (curve) {recovery$(L)$};
\draw[arr] (l3.east) -- (curve.west);
\end{tikzpicture}\\[2pt]
{\scriptsize swap one layer sim$\to$real; cascade vs.\ interior peak}
\end{minipage}
\caption{The GFIC protocol at a glance. Each axis compares a policy's internal
representation or action against an independent, causally-motivated reference (an
objective segmentation target for G, an input-level occlusion for F, a controlled
rendering pair for I, a paired activation swap for C) rather than a hand-labelled
semantic construct.}
\label{fig:gfic-overview}
\end{figure*}

\paragraph{G (Grounding).}
A linear probe is trained on the representation to recover object-ness segmentation,
using Segment Anything~\cite{kirillov2023sam} pseudo-ground-truth as an objective,
task-independent target (\cref{fig:gfic-overview}a). The primary readout is
\begin{equation}
G_m = \mathrm{mIoU}(\text{trained probe}) - \mathrm{mIoU}(\text{probe on random-init net}),
\end{equation}
a selectivity score against a control-task floor~\cite{hewitt2019designing}: the
random-initialization arm shares the probe's capacity and architecture but not its
learned features, so any gap above it reflects genuine content rather than probe
capacity. A second floor, \texttt{position\_only} (the probe given token coordinates
alone, no representation), bounds how much of the raw score is a spatial prior rather
than content; $G_m$ is a paired difference and remains valid even when that floor is
large.

\paragraph{F (Faithfulness).}
We test necessity directly at the input (\cref{fig:gfic-overview}b): the hazard entity is
occluded in the frame, and the resulting action is compared against the original
(hazard-present) action and the no-hazard baseline. With $a_{ghost}$, $a_{occ}$,
$a_{clean}$ the planned action under the three conditions, the primary readout is the
necessity ratio
\begin{equation}
R = \frac{a_{occ} - a_{clean}}{a_{ghost} - a_{clean}},
\end{equation}
$R \approx 0$ meaning the action fully reverts once the entity is removed (necessity
holds) and $R \approx 1$ meaning occlusion leaves the action unchanged from the
hazard-present case -- a blind-action signature -- i.e.\ the entity is not what drives it.
A grey-patch control arm, occluding an equal area elsewhere in the frame, rules out that
occlusion by itself (independent of location) moves the action.

\paragraph{I (Invariance).}
\Cref{fig:gfic-overview}c. Given a controlled appearance-only domain pairing
($do(\text{appearance})$, geometry and
ground truth held fixed), we report two quantities: the representational
\begin{equation}
D_L = \frac{\mathbb{E}\lVert Z_L(x_{real}) - Z_L(x_{sim})\rVert}{\mathbb{E}\lVert Z_L(x)\rVert},
\qquad I_m = 1 - D_{L^*},
\end{equation}
and the behavioural domain sensitivity
$\mathbb{E}|\Delta v_{cmd}| / \overline{v_{cmd}}$ on the same pairing. Invariance is not
reported as a single number: as \cref{sec:experiments} shows, the two readouts can order
the same models oppositely, which is itself diagnostic information a scalar would erase.

\paragraph{C (Concentration).}
\Cref{fig:gfic-overview}d. On the same domain pairing, we perform layer-wise activation
patching in which the
corruption is always a paired real-input swap (never synthetic noise) and the recovery
metric is always continuous (never binarized):
$\mathrm{recovery}(L) = $ the fraction of the sim-vs-real behavioural gap closed by
patching layer $L$'s activations from the sim run into the real run. The summary
statistic $C_m$ = top-2-layer share of total recovery is reported only after a shape
diagnostic confirms the recovery profile has an interior peak (\cref{sec:experiments});
when the profile is instead monotonic (a cascade), we report profile shape statistics
(Spearman correlation between layer index and recovery, responsible-layer entropy)
instead, since $C_m$'s premise is violated and its value is not meaningful.

\paragraph{Why no composite score.}
We deliberately do not combine G, F, I, C into a single weighted score. The four axes
differ in units, in adjudication state (some cells are indeterminate, some are not
comparable across a given pair of models), and in what a violated premise means for
their formula. Collapsing them into a scalar would require assigning weights with no
independent justification and would make the provenance of any resulting ranking
unauditable. We instead report the full matrix, with explicit ``not comparable'' and
``indeterminate'' labels treated as part of the result.
